\documentclass[12pt]{article}
\usepackage[T1]{fontenc}
\usepackage[utf8]{inputenc}
\usepackage{lmodern}
\usepackage{textcomp}
\usepackage{enumitem}
\usepackage{geometry}
\geometry{margin=1in}
\begin{document}
\begin{center}
Answer Key - Philosophy - Graduate Level Assessment 8 \
\small (Answers are ordered exactly as in the question paper.)
\end{center}

\section*{Multiple Choice}
\begin{enumerate}[label=\arabic*.]
\item \textbf{Answer:} A
\\ \textit{Options:} A) Babylonia and Europe; B) Rome and Europe; C) Palestine and France; D) Palestine and Rome; E) Europe and Egypt
\\ \textit{Note:} The correct answer identifies Babylonia and Europe as the two main centers for Jewish development post-Bar Kochba revolt because Babylonia became a hub for Jewish scholarship and religious authority, while Europe saw the emergence of significant Jewish communities and cultural contributions, especially in the medieval period.
\item \textbf{Answer:} A
\\ \textit{Options:} A) Sacrifices; B) Blessings; C) Offerings; D) Prayers; E) Apparitions
\\ \textit{Note:} Murtis, which are physical representations of deities in Hinduism, can be understood as sacrifices in the philosophical context of offering oneself and one's devotion to the divine through the act of worship and ritual.
\item \textbf{Answer:} A
\\ \textit{Options:} A) its rarity or frequency.; B) societal norms and values.; C) the individual's personal preference.; D) the amount of effort required to obtain it.; E) the potential pain that might accompany it.
\\ \textit{Note:} According to Mill, the value of a particular pleasure is influenced by its rarity or frequency because he argues that pleasures that are less commonly experienced tend to be regarded as more valuable due to their scarcity and the unique satisfaction they provide.
\item \textbf{Answer:} A
\\ \textit{Options:} A) must take into account the interests of all living beings.; B) make take into account the interests of all sentient beings.; C) should primarily focus on preserving the natural environment.; D) must align with societal norms and expectations.; E) are based solely on religious doctrines.
\\ \textit{Note:} Baier's ethical framework emphasizes the importance of considering the well-being and interests of all living beings in the formulation of genuine moral rules, advocating for a comprehensive moral perspective that includes all sentient entities.
\item \textbf{Answer:} D
\\ \textit{Options:} A) are not based on logical reasoning.; B) are flawed due to lack of empirical evidence.; C) contradict the principles of human rights.; D) follow common sense.; E) are based on biased principles.
\\ \textit{Note:} Reiman contends that van den Haag's arguments for the death penalty resonate with common sense because they appeal to intuitive moral reasoning and societal beliefs about justice and deterrence, which often prioritize retribution and the protection of society over nuanced ethical considerations.
\end{enumerate}

\section*{True / False}
\begin{enumerate}[label=\arabic*.]
\setcounter{enumi}{5}
\item \textbf{Answer:} False
\\ \textit{Options:} A) True; B) False
\\ \textit{Note:} Anscombe argues that moral psychology should focus on individual intention and action rather than societal norms, emphasizing the importance of personal virtue over external expectations.
\item \textbf{Answer:} False
\\ \textit{Options:} A) True; B) False
\\ \textit{Note:} Wolf argues that a moral saint prioritizes moral virtues over personal pleasures and leisure activities, suggesting that the ideal of moral sainthood does not necessarily include engaging in gourmet cooking and sports.
\item \textbf{Answer:} False
\\ \textit{Options:} A) True; B) False
\\ \textit{Note:} Abolishing the death penalty does not imply that all crimes are without consequences, as it is possible to advocate for alternative forms of punishment that address criminal behavior while rejecting capital punishment.
\item \textbf{Answer:} True
\\ \textit{Options:} A) True; B) False
\\ \textit{Note:} In a hypothetical syllogism, the structure is such that if the minor premise affirms the antecedent of a conditional statement, then it logically follows that the conclusion must affirm the consequent.
\item \textbf{Answer:} True
\\ \textit{Options:} A) True; B) False
\\ \textit{Note:} Hume argues that reason is subordinate to passion, asserting that while reason can inform and guide our actions, it is ultimately our emotions and desires that drive human behavior and emotional responses, making the statement true in the context of his philosophy.
\end{enumerate}

\section*{Long Form Response}
\begin{enumerate}[label=\arabic*.]
\setcounter{enumi}{10}
\item \textbf{Answer:} The externalist approach to meaning in philosophy posits that the significance of concepts and the understanding of language are shaped largely by factors external to the individual, such as social context and environmental factors. From an existentialist perspective, this externalism emphasizes that meaning is not an inherent quality of objects or experiences but is forged through individual engagement with the world. Existentialism asserts that existence precedes essence, suggesting that individuals must create their own meaning through lived experiences and choices. This notion challenges the idea of pre-determined meaning, encouraging a dynamic interaction between the self and the external world. Therefore, the existentialist externalist standpoint leads to a conception of meaning as fluid and contingent upon personal and societal contexts, influencing how individuals navigate existence and assign significance to their lives.
\\ \textit{Note:} correct because it accurately describes how the externalist approach highlights the role of social and environmental contexts in shaping meaning, while also aligning with existentialist ideas that prioritize individual experience and the notion that meaning is created through engagement with existence rather than being predetermined.
\item \textbf{Answer:} Pence argues that the ethical implications of somatic cell nuclear transfer (SCNT) must be evaluated in the context of its potential benefits, rather than solely on the basis of possible genetic abnormalities. He contends that the fear of genetic defects should not be a decisive objection to SCNT, as all medical procedures carry inherent risks. Moreover, advancements in genetic research and technology may mitigate these concerns, allowing for better screening and correction of abnormalities. By focusing on the potential for SCNT to alleviate suffering and advance medical science, Pence believes that the ethical discourse should prioritize the positive outcomes rather than dwelling on speculative risks. Thus, he positions the potential for genetic abnormalities as an insufficient basis for rejecting the practice of SCNT outright.
\\ \textit{Note:} The answer correctly highlights Pence's argument that the ethical evaluation of SCNT should focus on its potential benefits and the inherent risks of all medical procedures, while also acknowledging the advancements in genetic technology that could address concerns about genetic abnormalities.
\end{enumerate}

\vfill
\noindent\textit{End of Answer Key}
\end{document}